% Source: https://artofproblemsolving.com/community/c3h2364441_grrrrrrrrrrr
% Stable identity: aops-topic-2364441-q01
\textbf{Formula.} With $k=0$, and since $0!=1$,
\[\dbinom n0=\dfrac{n!}{0!\,n!}=\dfrac{n!}{1\cdot n!}=1.\]

\textbf{Counting.} $\dbinom n0$ is the number of ways to choose $0$ objects from $n$ objects. There is exactly one way: choose nothing (the empty set). So $\dbinom n0=\boxed{1}$ for every $n\ge0$, including $n=0$.
